\chapter{S}
\label{chap:s}

\term{S-matrix (Scattering Matrix)}
In quantum mechanics, a matrix that relates the initial and final states of a physical system that undergoes a scattering process. The elements represent the probability amplitudes for transitions between different states.

\term{S-adic Space}
A space constructed as the limit of a sequence of spaces, often used in the study of tilings and symbolic dynamics. It generalizes the notion of $p$-adic spaces to a sequence of different bases.

\term{Sample}
A subset of a population selected for statistical study; random samples aim to be representative.

\term{Sample Space}
The set of all possible outcomes of a random experiment, denoted $\Omega$.

\term{Scalar}
A quantity that is described by a single number, as opposed to a vector which has both magnitude and direction. In linear algebra, scalars are the elements of the field over which a vector space is defined.

\term{Scalar Curvature}
The trace of the Ricci curvature, $R = g^{ab}\operatorname{Ric}_{ab}$, a key invariant of Riemannian manifolds.

\term{Scalar Product}
Another name for the inner product. For two vectors $\mathbf{u}$ and $\mathbf{v}$ in $\mathbb{R}^n$, it is the sum $\sum u_i v_i$. It is used to define lengths and angles between vectors.

\term{Scale Factor}
The ratio by which the dimensions of a figure are multiplied under a similarity transformation.

\term{Scale Space}
A multi-scale representation of an image or signal, obtained by repeatedly smoothing the original data with a low-pass filter (typically Gaussian). It is used in computer vision to detect features at different resolutions.

\term{Schauder Basis}
A sequence of vectors $\{e_n\}$ in a Banach space $X$ such that every $x \in X$ has a unique representation $x = \sum_{n=1}^\infty a_n e_n$. Unlike Hamel bases, Schauder bases allow for infinite sums.

\term{Scheme}
The central object of study in modern algebraic geometry. A scheme is a locally ringed space that is locally isomorphic to the spectrum of a commutative ring. It generalizes the notion of an algebraic variety.

\term{Schröder--Bernstein Theorem}
A classical result in set theory, also called the Cantor--Bernstein theorem: if there are injections from $A$ into $B$ and from $B$ into $A$, then there exists a bijection between $A$ and $B$. It follows that cardinal comparison is antisymmetric, so two sets are equinumerous precisely when each is no larger than the other.

\term{Schur Complement}
For a block matrix $M = \begin{pmatrix} A & B \\ C & D \end{pmatrix}$, the Schur complement of the block $D$ is the matrix $A - B D^{-1} C$. It is used in matrix inversion and the study of positive-definite matrices.

\term{Schur's Lemma}
A fundamental result in representation theory: any irreducible representation of a group $G$ over an algebraically closed field has only scalar multiples of the identity as endomorphisms.

\term{Schur Decomposition}
The factorization of a square matrix $A$ as $A = QUQ^*$ where $Q$ is unitary and $U$ is upper triangular. The diagonal elements of $U$ are the eigenvalues of $A$.

\term{Schwarz Lemma}
A result in complex analysis: if $f$ is a holomorphic function from the unit disk to itself such that $f(0)=0$, then $|f(z)| \leq |z|$ and $|f'(0)| \leq 1$ for all $z$ in the disk.

\term{Schwarz Space}
The space of smooth functions whose derivatives all decay faster than any power of $1/|x|$ as $|x| \to \infty$. It is the natural domain for the Fourier transform, as the transform of a Schwarz function is also a Schwarz function.

\term{Schwartz Space}
The space $\mathcal{S}(\mathbb{R}^n)$ of rapidly decreasing smooth functions, on which the Fourier transform is an isomorphism.

\term{S-curve}
A sigmoid-shaped curve used in various fields to describe growth processes that start slowly, accelerate, and then level off as they reach a carrying capacity.

\term{Second Isomorphism Theorem}
If $G$ is a group with a subgroup $H$ and a normal subgroup $N$, then $HN/N \cong H/(H \cap N)$. This is one of the isomorphism theorems; the ring and module analogues state the same isomorphism when $N$ is an ideal or a submodule.

\term{Second-order PDE}
A partial differential equation where the highest derivative is of the second order. These are classified as elliptic, parabolic, or hyperbolic based on the coefficients of the second-order terms.

\term{Section (Fiber Bundle)}
A continuous map $s: B \to E$ from the base space $B$ to the total space $E$ of a fiber bundle such that $\pi \circ s = \text{id}_B$. A section picks one point from each fiber.

\term{Sector (Geometry)}
The portion of a disk enclosed by two radii and an arc. The area of a sector with radius $r$ and angle $\theta$ is $\frac{1}{2} r^2 \theta$.

\term{Seifert--van Kampen Theorem}
If $X = U \cup V$ with $U$, $V$, and $U \cap V$ path-connected, the fundamental group of $X$ is the amalgamated free product of $\pi_1(U)$ and $\pi_1(V)$ over $\pi_1(U \cap V)$. It allows the computation of fundamental groups by decomposing spaces into simpler parts: for instance, the wedge of $k$ circles has free group on $k$ generators, while $\pi_1(S^n)$ is trivial for $n \geq 2$.

\term{Self-adjoint Operator}
An operator equal to its adjoint, $T = T^*$; on Hilbert spaces self-adjoint operators have real spectrum and a spectral decomposition.

\term{Semi-algebraic Set}
A subset of $\mathbb{R}^n$ defined by a finite number of polynomial inequalities and equalities. They are the basic objects of study in real algebraic geometry.

\term{Semi-continuous Function}
A function is lower semi-continuous if its lower limit at any point is at least the function's value at that point; upper semi-continuous if the upper limit is at most the value.

\term{Semi-group}
A set equipped with an associative binary operation. Unlike a group, a semi-group does not require an identity element or inverses.

\term{Semi-norm}
A function that satisfies all the properties of a norm except for the identity of indiscernibles (i.e., $\|v\|=0$ does not imply $v=0$). It is used to define the topology of certain function spaces.

\term{Separable Space}
A topological space that contains a countable dense subset. For example, $\mathbb{R}$ is separable because $\mathbb{Q}$ is a countable dense subset.

\term{Separable Extension}
An algebraic extension $L/K$ where every element of $L$ is the root of a separable polynomial over $K$ (a polynomial with distinct roots in its splitting field).

\term{Separating}
(Topology) A family of functions or sets that distinguishes distinct points, as in the axioms of a Tychonoff space.

\term{Sequence}
An ordered list of elements, typically indexed by the natural numbers. A sequence $(a_n)$ converges to $L$ if for every $\epsilon > 0$ there exists $N$ such that $|a_n - L| < \epsilon$ for all $n \geq N$.

\term{Series}
The sum of the terms of a sequence. A series converges if the sequence of its partial sums has a finite limit.

\term{Set Theory}
The branch of mathematics studying sets and their operations, relations, and cardinalities, forming the foundation of mathematics.

\term{Sheaf}
A presheaf satisfying the locality and gluing axioms, allowing consistent local data to be assembled into global sections.

\term{Sieve Method}
A technique in number theory for counting or estimating the number of elements in a set that are not divisible by any of a given set of primes. The Sieve of Eratosthenes is the most basic example.

\term{Sigma-Algebra}
A collection of subsets of a set closed under complements and countable unions, providing the measurable sets of a measure space.

\term{Sigmoid Function}
A mathematical function having a characteristic "S"-shaped curve. The most common is the logistic function $f(x) = 1/(1+e^{-x})$, used in neural networks and population modeling.

\term{Simplex}
The generalization of a triangle to any dimension. An $n$-simplex is the convex hull of $n+1$ affinely independent points. A 0-simplex is a point, a 1-simplex is a line segment, and a 2-simplex is a triangle.

\term{Simple Group}
A group whose only normal subgroups are the identity and the group itself. Simple groups are the "building blocks" of all finite groups via composition series.

\term{Simple Loop}
A continuous path in a topological space that starts and ends at the same point and does not intersect itself except at the endpoints.

\term{Simple Root}
In the theory of Lie algebras, a root that cannot be written as a sum of two other positive roots. The set of simple roots forms a basis for the root system.

\term{Simpson's Paradox}
A trend that holds in each separate subgroup of data can be reversed when the subgroups are pooled together. It arises from lurking or confounding variables and unequal group sizes; the classic example is a treatment that looks better for every patient group yet worse overall. (Relates to the existing "Simpson's Rule" only by name.)

\term{Singular Matrix}
A square matrix that is not invertible, which occurs if and only if its determinant is zero. A singular matrix has a non-trivial kernel.

\term{Singular Value Decomposition (SVD)}
A factorization of any $m \times n$ matrix $A$ as $A = U \Sigma V^T$ where $U$ and $V$ are orthogonal matrices and $\Sigma$ is a diagonal matrix of singular values. It is used in data compression and noise reduction.

\term{Singular Point}
A point on a curve or surface where the derivative vanishes or is undefined. At such a point, the curve may have a cusp, a node, or a self-intersection.

\term{Singular Value}
The square roots of the eigenvalues of $A^* A$. They measure the "stretch" of the linear transformation $A$ along specific directions.

\term{Sine (Trigonometric)}
The ratio of the opposite side to the hypotenuse in a right triangle. For a unit circle, it is the $y$-coordinate of a point at angle $\theta$.

\term{Skew Lines}
Lines in space that are neither parallel nor intersecting.

\term{Skolem's Paradox}
By the Löwenheim--Skolem theorem, Zermelo--Fraenkel set theory (if consistent) has a countable model, yet it proves that uncountable sets exist. The resolution is that "uncountable" inside the model only means no internal bijection to the naturals; the apparent contradiction highlights the relativity of cardinality.

\term{Sleeping Beauty Problem}
A probability puzzle: Beauty is woken on Monday, and possibly again on Tuesday depending on a coin toss, without knowing which day it is. Whether the probability of heads is $1/2$ or $1/3$ is debated; the halfer and thirder positions rest on different interpretations of conditional probability and self-location.

\term{Smooth Manifold}
A topological manifold equipped with a smooth atlas of charts whose transition maps are infinitely differentiable.

\term{Sobolev Embedding Theorem}
Under suitable conditions on the indices, Sobolev spaces embed continuously into other Sobolev spaces, into $L^q$ spaces, or into spaces of continuous functions. For instance, $W^{k,p}$ embeds into $C^{m,\alpha}$ whenever $k - n/p > m$. Such embeddings play a central role in the regularity theory of PDEs, showing that weak solutions are smoother than initially guaranteed.

\term{Sobolev Solution}
A solution of a PDE that lies in a Sobolev space, defined through weak derivatives.

\term{Sobolev Space}
The space $W^{k,p}$ of functions whose weak derivatives up to order $k$ lie in $L^p$, the natural function space of PDE theory.

\term{Sorgenfrey Line}
The real line equipped with the lower-limit topology, where basis elements are half-open intervals $[a, b)$. It is a classic example of a space that is normal but not second-countable.

\term{Span (Linear Algebra)}
The set of all linear combinations of a given set of vectors. A set of vectors is a basis if it is linearly independent and its span is the entire vector space.

\term{Spectral Method}
A numerical method that expands the solution in global basis functions such as trigonometric or Chebyshev polynomials.

\term{Spectral Radius}
The maximum absolute value of the eigenvalues of a square matrix $A$. It determines the convergence of the power series $(I-A)^{-1} = \sum A^k$.

\term{Spectral Theorem}
A result stating that every symmetric matrix (or more generally, every normal operator on a Hilbert space) can be diagonalized by an orthonormal basis of eigenvectors.

\term{Spectral Theory}
The study of the spectrum (eigenvalues) of linear operators. It is fundamental to quantum mechanics, where observables are represented by self-adjoint operators.

\term{Sphere}
The set of all points in $\mathbb{R}^n$ at a fixed distance $r$ from a center point. An $n$-sphere $S^n$ is the boundary of an $(n+1)$-ball.

\term{Spherical Harmonic}
A set of special functions defined on the surface of a sphere. They are the angular portion of the solution to Laplace's equation in spherical coordinates.

\term{Spherical Geometry}
A non-Euclidean geometry where the "lines" are great circles on a sphere. In this geometry, the sum of the angles of a triangle is always greater than $180^\circ$.

\term{Spanning Tree}
A subgraph of a connected graph that is a tree and includes all the vertices of the original graph. Every connected graph has at least one spanning tree.

\term{Special Linear Group (SLn)}
The group of $n \times n$ matrices with determinant equal to 1. These transformations preserve volume and orientation.

\term{Special Orthogonal Group (SOn)}
The group of $n \times n$ orthogonal matrices with determinant 1. It consists of all rotations of $\mathbb{R}^n$ around the origin.

\term{Spec (Spectrum of a Ring)}
The set of all prime ideals of a commutative ring, equipped with the Zariski topology. It is the basic building block of schemes.

\term{Sperner's Theorem}
The largest antichain in the Boolean lattice of all subsets of an $n$-element set has size $\binom{n}{\lfloor n/2 \rfloor}$, realized by the middle level of subsets of size $\lfloor n/2 \rfloor$ or $\lceil n/2 \rceil$. It is a central result of extremal set theory and can be seen as a sharp version of Dilworth's theorem for the Boolean lattice.

\term{Sphere Packing Problem}
The problem of arranging non-overlapping spheres to occupy the maximum fraction of space. In 3D, the Kepler conjecture states the maximum density is $\pi/\sqrt{18}$.

\term{St. Petersburg Paradox}
A lottery in which the payoff doubles with each successive head and the expected value diverges to infinity, yet people will pay only a small sum to play. It motivated Bernoulli's introduction of diminishing marginal utility of money (utility of the logarithm), and connects to expected value versus expected utility.

\term{Stability}
The property of a solution or numerical scheme that small perturbations remain small; in ODE theory, solutions starting near an equilibrium remain near it for all time.

\term{Stable Marriage Problem}
An algorithmic problem of finding a matching between two sets of elements such that no pair of elements would both prefer each other over their current partners.

\term{Stable Range}
In algebraic K-theory, the range of dimensions where certain stability properties of the general linear group $\operatorname{GL}_n(R)$ hold.

\term{Standard Form}
A conventional way of writing a mathematical expression to make it easy to read or compare, such as the standard form of a linear equation $Ax + By = C$.

\term{State Space}
In dynamical systems and control theory, a space whose coordinates represent the state of the system at a given time.

\term{Stationary Process}
A stochastic process whose statistical properties (mean, variance, autocorrelation) do not change over time.

\term{Stationary Point}
A point where the derivative of a function is zero. It is a candidate for a local maximum, minimum, or saddle point.

\term{Statistics}
The science of collecting, analyzing, interpreting, and presenting empirical data to make inferences about a population.

\term{Steady State}
A solution of a dynamical system that does not change with time, such as a fixed point of an iteration.

\term{Stochastic Integral}
The integral of a stochastic process with respect to a semimartingale, defined in the sense of Itô or Stratonovich.

\term{Stochastic Process}
A collection of random variables indexed by time or space. Examples include Brownian motion and Poisson processes.

\term{Stokes' Theorem}
For a compact oriented manifold $M$ with boundary $\partial M$ and a differential form $\omega$, it states $\int_M d\omega = \int_{\partial M} \omega$, relating the integral of the exterior derivative over the manifold to the integral of the form over its boundary. Green's theorem, Gauss's divergence theorem, and the classical curl theorem are all special cases.

\term{Stone-Cech Compactification}
The largest compact Hausdorff space $\beta X$ that contains a given completely regular space $X$ as a dense subspace.

\term{Stone-Weierstrass Theorem}
A generalization of the Weierstrass approximation theorem: any continuous function on a compact Hausdorff space can be uniformly approximated by elements of a dense subalgebra.

\term{Stratification}
The decomposition of a space into a disjoint union of smooth manifolds (strata) of varying dimensions, used in singularity theory.

\term{Strong Induction}
A form of mathematical induction where the truth of a statement for $n$ is proved by assuming the statement is true for all $k < n$.

\term{Strong Topology}
In the context of dual spaces, the topology on $V^*$ induced by the norm of the original space $V$.

\term{Sturm-Liouville Problem}
The eigenvalue problem $-(p y')' + q y = \lambda w y$ with boundary conditions, yielding systems of orthogonal eigenfunctions.

\term{Subalgebra}
A subset of an algebra that is itself an algebra under the same operations.

\term{Subbasis}
A collection of sets whose finite intersections form a basis of a topology.

\term{Subgradient}
For a convex function, any vector $g$ with $f(y) \geq f(x) + g \cdot (y - x)$ for all $y$; a generalisation of the gradient to non-smooth functions.

\term{Subgroup}
A subset of a group that is itself a group under the same operation.

\term{Submanifold}
A subset of a manifold that is itself a manifold, with the subspace topology and a compatible smooth structure.

\term{Subspace}
A subset of a vector space that is closed under addition and scalar multiplication.

\term{Sub-additive Function}
A function $f$ such that $f(x+y) \leq f(x) + f(y)$ for all $x, y$ in its domain.

\term{Super-additive Function}
A function $f$ such that $f(x+y) \geq f(x) + f(y)$ for all $x, y$ in its domain.

\term{Super-symmetry}
A hypothesized symmetry in physics relating bosons (integer spin) and fermions (half-integer spin).

\term{Support (of a function)}
The closure of the set of points where a function is non-zero: $\operatorname{supp}(f) = \overline{\{x : f(x) \neq 0\}}$.

\term{Supremum}
The least upper bound of a set. If a set has a maximum, the supremum is that maximum.

\term{Surjection}
A function that hits every element of its codomain, also called onto.

\term{Symmetric}
A relation $R$ with $xRy \Rightarrow yRx$; a symmetric matrix satisfies $A^T = A$.

\term{Symmetric Group (Sn)}
The group of all permutations of $n$ elements. It has order $n!$ and is a central object in finite group theory.

\term{Symmetric Matrix}
A square matrix $A$ such that $A = A^T$. All eigenvalues of a real symmetric matrix are real, and it is always orthogonally diagonalizable.

\term{Symmetric Space}
A Riemannian manifold where for every point $p$, there exists an isometry that acts as a geodesic reflection through $p$.

\term{Symmetric Polynomial}
A polynomial $p(x_1, \ldots, x_n)$ that is invariant under any permutation of its variables.

\term{Symmetrization}
The process of making a function or operator symmetric by averaging it over a group action.

\term{Symplectic Geometry}
The study of manifolds equipped with a closed, non-degenerate 2-form $\omega$. It is the natural geometry of phase spaces in classical mechanics.

\term{Symplectic Form}
A closed, non-degenerate differential 2-form $\omega$ on a manifold. It allows for a canonical isomorphism between the tangent and cotangent bundles.

\term{Symplectic Manifold}
A smooth manifold equipped with a symplectic form. All symplectic manifolds are even-dimensional.

\term{Synge's Equation}
A differential equation used in the study of the stability of geodesics in general relativity.

\term{Syntactic Category}
A category whose objects are types and whose morphisms are provable functional dependencies in a logical theory.

\term{S-system}
A system of differential equations used in biochemical systems theory to model the interaction of components.

\term{Sylow Theorems}
A set of theorems describing the $p$-subgroups of a finite group $G$: Sylow $p$-subgroups exist, any two are conjugate, and the number $n_p$ of them divides $|G|$ and satisfies $n_p \equiv 1 \pmod p$. They are indispensable for classifying groups of small order and for analyzing the structure of finite groups generally.
